% Source: sources/returns/S6_04_MAPS_MANIFOLDS/sections/fourfolds_D2_remaining_even.tex:2–440
\paragraph{Theorem D.2.45 (Proved: the remaining minimal residual-four, eight and twenty types).}
Retain the primitive Mumford filling at every multiplicative site,
no additional regular logarithmic transformations, and the indicated
choice at the type-$II$ site in the following eleven patterns:
\[
\begin{array}{c|c|c|c}
 \mathrm{OS/profile}&(a,b)&\text{zero indices}&\text{type-}II\text{ choice}\\\hline
 45/II&(1,8)&2,3&\text{original or good log}\\
 55/II&(5,4)&3&\text{original or good log}\\
 55/\mathrm{default}&(1,20)&2,3,4&\text{none}\\
 55/II&(1,20)&3&\text{original or good log}.
\end{array}                                                    \tag{D45.1}
\]
Indices are zero-based. In an original-only pattern the ordinary
clutch is $e\gamma$ with $e=\pm1$. At a good $II$ logarithm allow
every order $m=6k$, $k>0$, every integer ordinary clutch $e$, and
every power vector
\[
 v=r\delta+n\ell,\qquad r\in\mathbb Z,
 \qquad n-me=\sigma\in\{1,-1\},                             \tag{D45.2}
\]
where $\ell$ is the actual primitive fixed lift with
$\gamma^*(\ell)=1$. These are exactly the simply connected data
in this one-log sector; they are free by D.2.30.
All the resulting manifolds have $H_2=H_3=\mathbb Z/b$ and
zero free third Betti number. Their full Zhubr invariants, in
the primitive fibre coordinate specified below, are
\[
 \begin{array}{c|c|c|c|c}
 b&P&\lambda&\mu&\text{oriented diffeomorphism type}\\\hline
 4&2&0&0&Z_{4,1}\\
 8&0&0&0&Z_{8,0}\\
 20&14&8&0&Z_{4,1}\#Z_{5,1}.
 \end{array}                                                 \tag{D45.3}
\]
Here $Z_{4,1},Z_{8,0}$ are the explicit framed two-sphere surgery
models of D.2.44, and $Z_{5,1}$ is the nonzero refined-class,
zero fundamental-class model of D.2.36. In particular D.2.44--45
identify all eighteen minimal patterns of residual order
$4,8,20$, including every simply connected finite multiplicity
and twist in (D45.2). No assertion about additional regular
logs or higher translation weights is included.

\emph{Proof: the actual differential and its reference frame.}
The homology and fundamental-group statements are D.2.26,
D.2.30 and D.2.41. In particular $H^2(X;\mathbb Z)=0$, so
$c_1(X)=0$ and $X$ is spin. The differential construction of
D.2.44 extends to the stated caps with two precise changes.
At an original type-$II$ elliptic cusp, the smooth total
surface has base map
\[
 q=y^2-x^3,\qquad
 \ker(df)=\mathcal O(2y,3x^2,0)\oplus\mathcal O(0,0,1).
                                                               \tag{D45.4}
\]
To see that this is the map, not just the singular zero set,
solve the Weierstrass equation for the base using its nonzero
linear base coefficient, split off its nondegenerate $y$
direction, and take the cube root of the remaining holomorphic
unit. The extra scalar coordinate is a local product coordinate.
The critical ideal $(x^2,y)$ has transverse length two.
Its literal critical cycle is therefore $2E_{II}$.
The normal comparison is the normalized unitary matrix
\[
 \frac1{\sqrt{4|y|^2+9|x|^4}}
 \begin{pmatrix}2y&-3\overline{x}^{2}\\3x^2&2\overline y\end{pmatrix}.
                                                               \tag{D45.5}
\]
Its degree is two: the holomorphic map $(x,y)\mapsto(2y,3x^2)$
has two positive local preimages of a small regular value.
Thus the relative obstruction, as well as the scheme cycle,
has coefficient two. The scalar $\delta$ is nonzero in the
second abstract kernel summand, including on $E_{II}$.

At a good logarithmic cap the total cyclic quotient is free,
so its good cover is etale over the smooth quotient total
space. On that cover $q=s^m$. If $D$ is the reduced multiple
fibre, remove the common divisorial factor of the differential
and use the target line
\[
 \mathcal L=f^*T\mathbb P^1\otimes
                   \mathcal O(-(m-1)D).                     \tag{D45.6}
\]
The corrected differential pulls back to $m\,ds$ and is
surjective. The division is deck-equivariant, and a unit
change $q=u s^m$ gives $mu\,ds+s\,du$, still surjective.
Consequently a good logarithmic cap has no corrected critical
cycle. Counting the uncorrected divisorial zero would give
the wrong relative bundle. Its vertical kernel is a subbundle
and has the descended nowhere-zero scalar section.

Write $C$ for the complete corrected critical cycle: coefficient
one on each nodal elliptic curve and each triangular double
curve, coefficient two on an original $II$ elliptic curve,
and no contribution at a good logarithm. The locally free
kernel $\mathcal V$ has the global nonzero abstract section
$\delta$. Both $\mathcal V/\mathcal O\delta$ and $\mathcal L$
are topologically trivial because $H^2(X;\mathbb Z)=0$.
The framed bundle $\mathcal V\oplus\mathcal L$ therefore
provides a specified full reference frame $\rho$ of $TX$
away from $C$. The local matrices in D.2.44 and (D45.5) give
the exact relative class
\[
                  \lambda(TX,\rho)=-\operatorname{PD}(C).
                                                               \tag{D45.7}
\]
This specifies local coefficients, not merely the absolute
identity $\lambda=-c_2$. In particular we do not assume
that $C$ is nullhomologous; it has order five when $b=20$.

\emph{A relative-characteristic formula retaining the rational link.}
We prove the needed general fact. Let $N\subset X\setminus C$
be an oriented normally framed three-manifold generating
$H_3(X)=\mathbb Z/b$, and suppose that the stable tangent
frame induced by $\rho$ on $N$ is an even multiple of
$\nu\in\Omega^{\rm fr}_3=\mathbb Z/24$.
Choose a four-chain $B$ with $\partial B=bN$, transverse to
the curve strata and avoiding their finitely many singular
points. Set
\[
                     \mathcal L(N,C)=\frac{B\cdot C}{b}
                             \in\mathbb Q.                  \tag{D45.8}
\]
This is the actual rational number for the marked cycles,
not a residue modulo integers. It is independent of $B$,
since $[C]$ is torsion. Then
\[
       \operatorname{lk}(N,\operatorname{PD}P)
                          =-\frac12\mathcal L(N,C)
                             \quad\text{in }\mathbb Q/\mathbb Z.
                                                               \tag{D45.9}
\]

Here is a proof that includes the missing parity. Remove a
point away from all chains. Write $b=2^a q$ with $q$ odd.
Join $q$ parallel copies of $N$ by normally framed tubes.
Kaplan's prescribed-spin construction, followed by deletion
of its zero-handle and reversal, gives a cobordism to a
sphere using only two-handles
\cite[Definition 1.7, Remark 1.8 and Theorem 3.1]{mapsKaplan79}.
As in D.2.44, its attaching disks and normal four-frames can
be realized in the punctured simply connected six-manifold:
general position gives disjoint disks and the spin condition
is exactly their $\pi_1SO(4)$ extension obstruction.
Thus the outgoing sphere is actually embedded and represents
$q[N]$, the two-primary generator.
Now choose Zhubr's admissible three-frame $v$ for this
particular embedded sphere, as permitted by his Theorem 1
\cite{ffZhubr13}, and put $\xi=v^\perp$. Its spin framing
obstruction is $o_4(\xi)=P$, and
$\lambda(\xi)=2o_4(\xi)$ with his normalization.
Choose a spin frame $\tau$ of $\xi|_N$.
It extends over the tube and two-handle trace: the relative
obstruction groups are $\pi_0\mathrm{Spin}(3)$ and
$\pi_1\mathrm{Spin}(3)$, both zero. Hence the framings
induced by $v+\tau$ on the incoming and outgoing manifolds
are stably bordant. On the outgoing sphere, admissibility
means even framed parity, as normalized in D.2.44 and
\cite[Section 5]{ffZhubr13}. Since $q$ is odd,
$[N]_{v+\tau}$ is even as well.

The complement of the real codimension-four curve strata
is simply connected, by general position applied to disks.
Both $\rho|_N$ and $v+\tau$ therefore induce the same spin
structure. Their difference lifts to $\mathrm{Spin}(6)$.
Let $a$ be its normalized degree-three integral class and
$k=\langle a,[N]\rangle$. The stable $J$ image of the
$\pi_3\mathrm{Spin}(6)$ generator is $\nu$, so the two
even framed parities imply $k$ is even. The integral
transgression of $\lambda$ gives, with consistent
transition-map direction, the identity of relative classes
\[
 \lambda(TX,\rho|_N)=2o_4(\xi,\tau)+\delta a
                         \quad\text{in }H^4(X\setminus\{p\},N;\mathbb Z).
                                                               \tag{D45.10}
\]
If $Q=\langle o_4(\xi,\tau),B\rangle$, evaluation gives
the identity of integers
\[
                       -B\cdot C=2Q+bk.                    \tag{D45.11}
\]
By the integral definition of the torsion pairing,
$Q/b\bmod\mathbb Z=\operatorname{lk}(N,\operatorname{PD}P)$.
Dividing (D45.11) and using the evenness of $k$ proves
(D45.9). This is not division of a class in an even cyclic
group. Its integer correction was fixed by actual spin
frames before division. The proof uses Zhubr's proved
Theorem 1, and does not use his conjectural general
change-of-three-frame formula. It also permits a nonzero
odd-primary $\lambda$.

\emph{The finite-core generator and its reference parity.}
For a single good $II$ logarithm the regular three-torus
can be an even class in $H_3(X)$. We use the following
actual core three-manifold instead. The central four-manifold
is an elliptic-torus bundle over the rank-two lattice
\[
 \langle\delta,\ell,g\mid mg=r\delta+n\ell\rangle.
                                                               \tag{D45.12}
\]
Its elliptic monodromy is $B$ on $g$, trivial on the other
two generators, where $B$ has order six and
$\det(B-I)=1$. With $\sigma=n-me$, define the primitive
integer character
\[
 \chi(\delta)=1,\qquad \chi(\ell)=-\sigma r,
                       \chi(g)=-\sigma re.                 \tag{D45.13}
\]
The relation in (D45.12) is respected. Let $N$ be the
connected level set of the resulting map of the core to
$S^1$, cooriented by $\chi$. Put
$h=g+\sigma re\delta$ and $k_0=\ell+\sigma r\delta$.
These lie in $\ker\chi$ and satisfy $mh=nk_0$.
Since $\gcd(m,n)=1$, their primitive generator has elliptic
monodromy $B^\sigma$: an integer inverse of $n$ modulo
$m$ is congruent to $\sigma$ modulo six. Thus $N$ is an
actual affine elliptic mapping torus, with
\[
                        H_1(N;\mathbb Z)=\mathbb Z.        \tag{D45.14}
\]
Translations do not change this Wang sequence.

For completeness, any actual tangent parallelization of
a closed oriented three-manifold $M$ has framed parity
\[
        [M,\tau]_{\rm fr}\bmod2
                    =1+\dim H_1(M;\mathbb F_2).             \tag{D45.15}
\]
Use Kaplan's spin-compatible filling $W$ with one zero-handle
and $j$ two-handles. Its tangent bundle is trivial, since
its spine is two-dimensional and $w_2=0$. On the boundary
use the outward normal followed by $\tau$.
The relative stable framing obstruction is $\lambda$;
its $J$ image is $[M,\tau]$. Modulo two it is the relative
$w_4$, hence the relative Euler class, whose evaluation is
$\chi(W)=1+j$. The mod-two intersection matrix of the
spin filling is alternating, so has even rank. Its
cokernel is $H_1(M;\mathbb F_2)$; hence
$j\equiv\dim H_1(M;\mathbb F_2)$. This proves (D45.15)
while retaining the prescribed spin structure.

The reference frame on the finite-core $N$ is the
stabilization of an actual tangent frame. Indeed $\delta$
is transverse to $N$ by $\chi(\delta)=1$. For the global
complex vertical frame $(\delta,v_2)$, project
$J\delta,\operatorname{Re}v_2,\operatorname{Im}v_2$ to $TN$
along $\delta$. They form a tangent frame. The normal
frame is $\delta$ together with the two real directions
of the corrected target line $\mathcal L$. Projection
and horizontal splitting are upper triangular homotopies.
Equation (D45.14), then (D45.15), therefore gives even
reference parity. This conclusion is not inferred from
the even-degree torus cover of $N$.

\emph{The actual marked cycles and the unbounded chain identity.}
Use the same integral regular-fibre coordinate as D.2.41:
\[
                \Omega=a e_1^*e_2^*-b\delta^*\gamma^*.
                                                               \tag{D45.16}
\]
The full degree-two monodromy cokernel is $\mathbb Z$,
primitively measured by $\Omega$. Choose an integral
fibre two-cycle $s$ with $\Omega(s)=1$; its image generates
$H_2(X)=\mathbb Z/b$. The nodal and triangular preimages
are those proved in D.2.44. At an original $II$ site put
\[
 z=(B-I)^{-1}q_0\in\mathbb Z^2,\qquad
 \ell=(-z,0,1),\qquad [E_{II}]=\delta\wedge\ell,             \tag{D45.17}
\]
where $q_0$ is the full elliptic long-column vector.
The full matrix fixes $\ell$, not just its quotient.
The principal $\delta,\ell$ torus through the elliptic
cusp is the reduced critical elliptic curve; its regular
torus specializes with degree one. The comparison takes
place within that cap. Thus $2\delta\wedge\ell$ is the
literal critical preimage, with the complex torus orientation.
For OS45/II, $\ell=(-8,0,0,1)$; for the OS55/II rows,
$\ell=(0,0,0,1)$.

Transport all these preimages along the fixed base paths
and let $K=\Omega(C)$. The exact integers are
\[
\begin{array}{c|c|c|r|c}
 \text{profile}&(a,b)&\text{zero and }II\text{ choice}&K&
                                  \operatorname{ord}[C]\\\hline
 45/II&(1,8)&2,3;\ \mathrm{original}&-32&1\\
 45/II&(1,8)&2,3;\ \mathrm{log}&-16&1\\
 55/II&(5,4)&3;\ \mathrm{original}&-108&1\\
 55/II&(5,4)&3;\ \mathrm{log}&-100&1\\
 55/II&(1,20)&3;\ \mathrm{original}&-68&5\\
 55/II&(1,20)&3;\ \mathrm{log}&-28&5\\
 55/\mathrm{default}&(1,20)&2&-108&5\\
 55/\mathrm{default}&(1,20)&3,4&-68&5.
\end{array}                                                    \tag{D45.18}
\]
The difference $C-Ks$ bounds a chain in the regular
complement, after the specified local specialization
chains are included. This uses the full integral
monodromy image, whose five nonzero Smith factors are
units; it does not replace an image by its saturation.
For a logarithmic pattern all these chains avoid its
finite core, since that cap has no critical curve.

Here is the exact all-parameter relation that fixes
both the generator and the rational linking.
Every monodromy fixes $\delta$. Its quadratic bar
correction on that single loop is zero, including
the transported last peripheral word, and the
peripheral graph maps satisfy $\sum_iJ_i\delta=0$.
Take the finite normal Thom column for $\delta$,
add $m$ times the ordinary-clutch $\delta$ column,
and add $m$ times the other filled-cap $\delta$ sweeps.
The graph terms cancel. Its actual fibre part is
\[
 -\delta\wedge(r\delta+n\ell)+me\delta\wedge\gamma,
 \qquad \Omega=b(n-me)=b\sigma.                            \tag{D45.19}
\]
The remaining difference from $b\sigma s$ is an
integral monodromy relation, supported away from $N$.
All other cap sweeps also bound away from $N$.
The finite normal Thom chain is the disc bundle
over its actual central loop. Its intersection
with the level set $N$ is its character value,
which here is $\chi(\delta)=1$.
Thus (D45.19) gives the actual pairings
\[
 \operatorname{lk}(N,s)=\frac{\sigma}{b}\pmod{\mathbb Z},
 \qquad \mathcal L(N,C)=\frac{\sigma K}{b}.                 \tag{D45.20}
\]
The first is a unit pairing, proving that $N$ is a
full $H_3$ generator, including when $b=20$.
The fifth affine-meridian column, if used in another
Smith bounding chain, pairs by its actual value
$\chi(g)=-\sigma re$, not zero. This retained term
makes the full-chain comparison agree with (D45.20).
There is no additional factor of $m$: the covering
order is already present in the boundary norm column.

For an original-only pattern use the ordinary
$T^3=\langle e_1,e_2,\gamma\rangle$ and $e=\pm1$.
The same peripheral relation, now using its ordinary
normal Thom chain, has fibre value $-be$ and
intersection one with $T^3$. Therefore
\[
 \operatorname{lk}(T^3,s)=-\frac e b\pmod{\mathbb Z},
 \qquad \mathcal L(T^3,C)=-\frac{eK}{b}.                    \tag{D45.21}
\]
It is again a full generator with even reference parity,
as proved in D.2.44. If $b=20$, these equations can
equivalently be computed by bounding $5C$ and dividing
its literal intersection by five. No nullhomology
of $C$ is asserted.

Substitute (D45.20) or (D45.21) into the already
proved relative formula (D45.9). Since $K$ is an
even integer in every row of (D45.18), cancellation
of the displayed unit pairing gives
\[
                       P=-K/2\quad\text{in }\mathbb Z/b.
                                                               \tag{D45.22}
\]
Here the division was justified by (D45.10)--(D45.11)
using the actual cycles and framings. It is not
division of the congruence $2P=\lambda$.
The table gives $P=2,0,14$ as in (D45.3), and the
absolute identity $\lambda=-[C]$ gives $0,0,8$.

\emph{The complete fundamental-class invariant and recognition.}
The global auxiliary-line comparison and localization
of D.2.44 apply over $\mathbb Z/4$ and $\mathbb Z/8$,
respectively. All actual original-$II$ coefficients
are zero and all finite-$II$ normal coefficients
are $bn$ and $-br$; hence the same mod-$b$ class
extends for every parameter in (D45.2).
For $b=20$ apply that calculation separately over
$\mathbb Z/4$ and $\mathbb Z/5$. The $C_4$ action
has only the rational fixed curve and the $I_4$
elliptic fixed curve, with characters $2,0$.
The $C_5$ action similarly has the rational and
$I_5$ curves, with characters $0,0$; all remaining
stabilizers are prime to the relevant coefficient
order. The actual Tate normal weights have square
one modulo four, and weight two modulo five.
The cube formula is zero in each coefficient ring:
on the rational curve it is $-3$ times the square
of its point character, and on the elliptic curve
the possible mixed degree-one term has square zero.
Thus the full fundamental-class invariant is zero,
since $H_6(K(\mathbb Z/b,2);\mathbb Z)=\mathbb Z/b$
is detected by the cube for $b=4,8,20$
\cite[Sections 3.11--3.14]{manZhubr75}.

The finite data therefore agree with the first two
surgery models of D.2.44. For $b=20$, the two
primary parts of $P=14$ are $2\in\mathbb Z/4$
and a nonzero element of $\mathbb Z/5$.
The inclusion
$K(\mathbb Z/4,2)\vee K(\mathbb Z/5,2)
\to K(\mathbb Z/4,2)\times K(\mathbb Z/5,2)$
is a homology equivalence: all cross tensor and
Tor terms vanish because their positive-degree
homology has coprime primary orders. Both spaces
are simply connected, so it is a homotopy equivalence.
The reduced spin-bordism invariant consequently
splits into these two summands, as in D.2.38.
The complete invariants agree with
$Z_{4,1}\#Z_{5,1}$, and the free third Betti numbers
also agree. Zhubr's Theorems 3 and 5
\cite{manZhubr75} give all the oriented
diffeomorphisms in (D45.3).

\emph{Exact computations and their bounds.}
The local differential identities are checked in
\path{code/maps_D2_II_frame_checks.py}.
The line/support, full angular-image, actual
adapted-matrix and arbitrary finite-normal
coefficients for all eighteen patterns are checked in
\path{code/fourfolds_D2_remaining_AH.py}.
This comprises fourteen initial rows, $1056$ exact
semicharacter checks, $62$ full adapted matrices,
four symbolic arbitrary-twist finite normal rows
and four original-$II$ choices.
The actual curve vectors, complete Smith solutions,
finite root-loop characters and literal intersections
are retained by
\path{code/manifolds_extended_critical_cycles.py}.
Its $158$ full Smith replays comprise the two original
clutch signs and, for each of the four logarithmic
patterns, $k=1,2,3$, $r=0,1$, $e=-1,0,1$ and
$\sigma=\pm1$. This explicitly bounded replay is
supplemented by the eleven constant certificates in
\path{records/manifolds_extended_critical_symbolic.json}:
integral preimages of $-\delta\ell-bs$,
$\delta\gamma+bs$, and $C-Ks$, together with the
zero bar corrections and exact peripheral graph
identity. Their linear combinations prove (D45.19)
for every integer parameter in (D45.2), with no
unproved finite search bound.\hfill$\square$
